% =====================================================================
%  Paper I  --  KHI en RRMHD  --  ESQUELETO MNRAS
%  Clase y .bst tomados de CTAN (macros/latex/contrib/mnras), v3.2.
%  Guia completa de la clase: mnras_guide.pdf (en esta misma carpeta).
%
%  COMO COMPILAR:  ./build.sh        (pdflatex + bibtex + pdflatex x2)
%
%  ANTES DE ESCRIBIR, LEER:
%    ../manuales/MNRAS_instrucciones.md   requisitos de la revista
%    ../02_PLAN_EXTRACCION.md             que sale de donde en la monografia
%    ../04_ESTILO_INGLES.md               ingles britanico + glosario ES->EN
%    ../05_CHECKLIST_ENVIO.md             checklist antes de enviar
% =====================================================================

% MNRAS pide \pdfminorversion=5 con pdfTeX, para ScholarOne.
\pdfminorversion=5

\documentclass[fleqn,usenatbib]{mnras}

\usepackage{newtxtext,newtxmath}   % MNRAS se compone en Times
\usepackage[T1]{fontenc}
\DeclareRobustCommand{\VAN}[3]{#2}
\let\VANthebibliography\thebibliography
\def\thebibliography{\DeclareRobustCommand{\VAN}[3]{##3}\VANthebibliography}

\usepackage{graphicx}
\usepackage{amsmath}
% OJO: no cargar amssymb junto con newtxmath (conflicto documentado por MNRAS).

% Las figuras se toman directamente de la monografia mientras no se hayan
% regenerado en ingles y al ancho de columna correcto. Ver ../03_INVENTARIO_FIGURAS.md
\graphicspath{{figures/}{../../monografia/figuras/}}

% --- Macros del proyecto (minimas, en notacion MNRAS) --------------------
% Subindices-etiqueta en REDONDA (\mathrm), subindices-variable en cursiva.
\newcommand{\gKHI}{\gamma_{\mathrm{KHI}}}
\newcommand{\scrit}{\sigma_{\mathrm{crit}}}
\newcommand{\Rmst}{\mathrm{Rm}^{*}}
\newcommand{\Ozp}{\Omega_{zp}}
\newcommand{\Omaxzp}{\Omega_{zp}^{\max}}
\newcommand{\Otot}{\Omega_{\mathrm{tot}}}
\newcommand{\akh}{a_{\mathrm{kh}}}
\newcommand{\vsh}{v_{\mathrm{sh}}}
\newcommand{\B}{\mathbfit{B}}      % vectores: negrita cursiva (MNRAS)
\newcommand{\Efield}{\mathbfit{E}}
\newcommand{\vel}{\mathbfit{v}}
\newcommand{\kvec}{\mathbfit{k}}
% Marcador de texto pendiente: \todo{...} sale en el PDF y es facil de buscar.
\newcommand{\todo}[1]{\textbf{[TODO: #1]}}

% =====================================================================
\title[Resistive KHI in RRMHD]{%
  \todo{titulo completo}%
  % Propuesta: The resistive-to-ideal transition of the relativistic
  % Kelvin--Helmholtz instability
}

% Encabezado corto: hasta 45 caracteres. Autores: et al. a partir de 3.
\author[B. Martinez Anzola et al.]{%
  B. Martinez Anzola,$^{1}$\thanks{E-mail: \todo{correo}}
  S. Rodriguez Garcia,$^{1}$
  and S. Miranda-Aranguren$^{1}$
  \\
  % \newauthor  % <- usar si hace falta otra linea de autores
  $^{1}$Programa Academico de Fisica, Facultad de Ciencias Matematicas y
  Naturales, Universidad Distrital Francisco Jose de Caldas,
  Bogota, Colombia\\
  % TODO: confirmar afiliaciones exactas y orden de autoria (ver ../05_CHECKLIST_ENVIO.md)
}

\date{\todo{fecha}}
\pubyear{2026}

\begin{document}
\label{firstpage}
\pagerange{\pageref{firstpage}--\pageref{lastpage}}
\maketitle

% =====================================================================
\begin{abstract}
\todo{Maximo 250 palabras, UN SOLO parrafo. Debe contener: (i) que se hizo
(barrido de 44 simulaciones RRMHD en $\sigma\in[100,10500]$ con el codigo
\textit{Cueva}); (ii) el modelo sigmoide con \emph{offset} y la asintota ideal
$C+\gamma_0=1.03\pm0.05$, consistente al 8 per cent con la teoria lineal RMHD
compresible; (iii) que la transicion resistivo-ideal es una ZONA
$\sigma\in[1400,7000]$ con centro $\scrit\approx2815$; (iv) la ley
$\Omaxzp\propto\sigma^{1.22}$, por encima del piso de Sweet--Parker
$\sigma^{1/2}$; (v) una frase sobre la campana magnetica. Escribirlo AL FINAL.
Base en espanol: monografia/sections/00\_preliminares.tex}
\end{abstract}

% Entre 1 y 6, de la lista oficial de MNRAS (verificar el PDF de keywords).
\begin{keywords}
instabilities -- magnetic reconnection -- MHD -- plasmas --
relativistic processes -- methods: numerical
\end{keywords}

% =====================================================================
\section{Introduction}
\label{sec:intro}

\todo{~1.5 columnas. Fuente: monografia/sections/01\_introduction.tex.
Cinco parrafos: (1) KHI en jets relativistas y AGN; (2) por que resistiva y no
ideal -- el congelamiento del flujo prohibe la reconexion (ESTE ES EL GANCHO);
(3) trabajo previo -- REQUIERE BUSQUEDA BIBLIOGRAFICA NUEVA, ver
../02\_PLAN\_EXTRACCION.md seccion 7.1; (4) que hace este trabajo, como
preguntas de investigacion y NO como "objetivos especificos"; (5) estructura
del articulo.}

% =====================================================================
\section{Numerical method and setup}
\label{sec:method}

\subsection{Governing equations}
\label{sec:equations}

\todo{Media a una columna. Solo el sistema RRMHD en forma conservativa 3+1 y la
ley de Ohm relativista, citando Komissarov (2007); una frase sobre el sistema
aumentado GLM con $\Psi$ y $\Phi$. Fuente: 02\_rrmhd\_theory.tex l. 373+ y 413+.
NO reproducir la derivacion del tensor de Faraday ni los lagrangianos.}

\subsection{The \textit{Cueva} code}
\label{sec:code}

\todo{Un parrafo denso: volumenes finitos, MP5, HLL, GLM, IMEX-RK. Mas SOLO lo
no estandar: condicion subcaracteristica e inversion algebraica local del
termino de conduccion (04\_numerical\_methods.tex l. 75--81 y 148--168).
Confirmar con el director la cita correcta del codigo y si hay DOI/ASCL.}

\subsection{Initial conditions}
\label{sec:ic}

\todo{Doble capa de cizalla tipo Mizuno et al. (2013), test 35A:
$\vsh=0.5c$, $\akh=0.05$, $\B=(0,\sqrt{0.02},\sqrt{2.0})$, $\Gamma=4/3$,
dominio 2D, malla $512\times256$. Fuente: setup\_exp.tex l. 24--164.
La Tabla~\ref{tab:setup} sale de tab:setup\_summary.}

\begin{table}
  \centering
  \caption{Parameters of the reference setup.
           \todo{traducir tab:setup\_summary de setup\_exp.tex:139}}
  \label{tab:setup}
  \begin{tabular}{lll}
    \hline
    Quantity & Symbol & Value \\
    \hline
    Shear velocity     & $\vsh$  & $0.5\,c$ \\
    Shear width        & $\akh$  & $0.05$ \\
    Adiabatic index    & $\Gamma$ & $4/3$ \\
    Grid               &          & $512\times256$ \\
    \todo{completar}   &          & \\
    \hline
  \end{tabular}
  % MNRAS: SOLO lineas horizontales, nunca verticales.
\end{table}

\subsection{The two campaigns}
\label{sec:campaigns}

\todo{(a) Barrido: 44 simulaciones, $\sigma\in[100,10500]$; explicar que
$\sigma=10500$ fue el LIMITE NUMERICO del codigo (o truncar a $10^4$ y
justificarlo). La tabla canonica completa (44 filas) va a material
suplementario, no al cuerpo -- MNRAS pide 5--10 filas de muestra.
(b) Campana magnetica: fundir tab:campana\_A/B/C en una sola tabla.}

% =====================================================================
\section{Results: the conductivity sweep}
\label{sec:results}

\subsection{Global behaviour of the sweep}
\label{sec:overview}
\todo{Fuente: 05\_results\_discussion.tex sec:vision\_global.
Figuras: fig:ultrawide\_omega\_zp, fig:mapa\_conductividad.}

\subsection{Extraction of the linear growth rate}
\label{sec:fit}
\todo{Fuente: sec:lnfit. Figura fig:lnfit (quiza a apendice).}

\subsection{The sigmoid model and the ideal asymptote}
\label{sec:sigmoid}
\todo{RESULTADO CENTRAL. Fuente: sec:sigmoide\_v6 + tab:par\_v6.
Figura fig:ajuste\_v6 = LA figura del articulo.
Declarar explicitamente que $C+\gamma_0$ es una EXTRAPOLACION: el dato de mayor
$\sigma$ alcanza el 85 per cent del plateau.}

\subsection{Statistical validation}
\label{sec:validation}
\todo{Comprimir mucho: AIC para la seleccion de modelo, jackknife para las
incertidumbres, hold-out para la validacion. Una o dos frases; el detalle a
apendice o suprimido (es estandar). Fuente: sec:validacion\_v6.}

\subsection{The transition is a zone, not a threshold}
\label{sec:zone}
\todo{SEGUNDO RESULTADO. $\sigma\in[1400,7000]$, centro $\scrit\approx2815$;
cada observable se idealiza a una $\sigma$ distinta. Fuente: sec:scrit\_v6,
figura fig:zona\_transicion, tabla tab:est\_v6.}

\subsection{Power-law scaling of the peak enstrophy}
\label{sec:powerlaw}
\todo{TERCER RESULTADO. $\Omaxzp\propto\sigma^{1.22}$, por encima del piso
Sweet--Parker $\sigma^{1/2}$ -> firma de reconexion, con \emph{tearing}
secundario como HIPOTESIS (no se mapearon las laminas de corriente).
Fuente: sec:leyes\_potencia, fig:omega\_max\_pot, tab:leyes\_potencia.}

\subsection{Global electromagnetic variables}
\label{sec:global}
\todo{$J_{\max}$ crece con $\sigma$; el trabajo EM acumulado aumenta; la
amplificacion de energia magnetica cae de ~15 per cent (resistivo) a nula
(ideal). $f_{Vz}$ maxima en la zona de transicion. Fuente: sec:globales.}

% =====================================================================
\section{Comparison with linear theory}
\label{sec:theory}

\todo{Techo inviscido (Rayleigh) -> 1 parrafo. Prediccion RMHD compresible
(Lees--Lin) con $c_s\to v_{f\perp}$ segun Chow et al. (2023):
$\hat\omega=0.095$ frente a nuestro $0.103$, 8 per cent.
$M_{re}=0.60<\sqrt2$ -> holgadamente inestable.
Fuente: sec:teoria\_v6 + 07\_appendices.tex seccion 1.}

% =====================================================================
\section{Magnetic field strength and orientation}
\label{sec:magnetic}

\todo{1.5 columnas para 25 pp de monografia. (a) Intensidad: $\gKHI$ cae de
0.81 a 0.38 y la enstrofia pico de ~87 a ~25 ($\sigma=10^4$); el agente es la
PRESION magnetica ($\beta\propto B_0^{-2}$), no la tension -- el flujo sigue
super-Alfvenico a lo largo de $\kvec$ ($M_{A,\parallel}=9.4$). (b) Orientacion:
solo $B_y$ (paralela a $\kvec$) ejerce tension, y solo a angulos grandes; $B_x$
cumple $\kvec\cdot\B=0$ y NO ejerce tension. (c) Balance energetico: antifase
$E_{\mathrm{kin}}\leftrightarrow E_{\mathrm{mag}}$; cuasi-reversible en
$\sigma=10^4$, amortiguado por disipacion ohmica en $\sigma=6000$.
Fuente: sec:campana\_magnetica y sec:balance\_energetico.}

% =====================================================================
\section{Discussion and conclusions}
\label{sec:conclusions}

\todo{Fuente: 06\_conclusions.tex, que ya esta en el registro correcto:
traducir y convertir la estructura "(1) objetivo... (2) objetivo..." en prosa
tematica. Conservar INTEGRAS las limitaciones (2D; EoS politropica; $\sigma$
escalar sin Hall; resolucion de las hojas resistivas; asintota extrapolada al
85 per cent). Trabajo futuro en 2--3 frases.}

% =====================================================================
\section*{Acknowledgements}

\todo{Recursos de computo (fuente: sections/agradecimientos\_cluster.tex,
comprimido a 2 frases), financiacion, personas, y el software usado
(numpy, scipy, matplotlib -- CITADOS en las referencias, no solo nombrados).
Declaracion de uso de IA: ver ../comun/DECLARACION\_IA.md.
Declaracion de contribucion estudiantil (premio MNRAS): ver ../05\_CHECKLIST\_ENVIO.md.}

% =====================================================================
% OBLIGATORIO en MNRAS, siempre despues de Acknowledgements.
\section*{Data Availability}

\todo{Ver ../comun/DATA\_AVAILABILITY.md -- usar la plantilla A (datos derivados
y scripts en Zenodo con DOI; datos crudos bajo peticion; \textit{Cueva} segun
decida el director).}

% =====================================================================
\bibliographystyle{mnras}
\bibliography{references}
% OJO: references.bib es copia de la monografia (biblatex). Ver
% ../comun/NOTAS_BIBTEX.md -- la recomendacion es reconstruirlo desde NASA ADS.

% =====================================================================
\appendix

\section{Supplementary statistical detail}
\label{app:stats}
\todo{Solo si hace falta. Fuente: 07\_appendices.tex seccion 3.
Casi todo (R\textsuperscript{2}, AIC, jackknife, hold-out, test de rachas) es
estandar y se resuelve con una frase en el cuerpo.}

\bsp	% typesetting comment
\label{lastpage}
\end{document}
